\documentclass{rapportECL}
\usepackage[utf8]{inputenc}
\usepackage[T1]{fontenc}
\usepackage[french]{babel}
% Gestion des marges pour A4 vertical
\usepackage[top=2.5cm, bottom=2.5cm, left=2.5cm, right=2.5cm]{geometry}
\usepackage{graphicx}
\usepackage{xcolor}
\usepackage{tikz}
\usetikzlibrary{arrows.meta, positioning, calc}
% Définition des styles pour les blocs, sommes, etc.
\tikzset{
    block/.style = {draw, rectangle, minimum height=2.5em, minimum width=4em, align=center, thick},
    sum/.style = {draw, circle, inner sep=2pt, thick},
    input/.style = {coordinate},
    output/.style = {coordinate},
    line/.style = {draw, -Latex, thick}
}
% Informations du rapport
\titre{Rendu Jumeaux numériques}\\
% \large UE : Dimensionnement des structures (ENISE-5GM)}
\eleves{Kevin \textsc{TONGUE}}
\enseignant{Enseignant : Renaud \textsc{FERRIER}}
\begin{document}
%----------- Initialisation -------------------
       
\fairemarges %Afficher les marges
\fairepagedegarde %Créer la page de garde
% \tabledematieres %Créer la table de matières

\centerline{22 d\'ecembre 2025}

\newpage

\section{4.1. Travail pr\'eparatoire}

\subsection{4.1.1. \'Etude du syst\`eme et de son asservissement}

Le syst\`eme est compos\'e d'un moteur pilot\'e en courant $i$, avec une fonction de transfert $F_m$ donnant le couple $C_m = F_m i$. La transmission a une fonction de transfert $F_t$ telle que la force $f_A = F_t C_m$. La poutre est repr\'esent\'ee par une fonction de transfert $F$ ayant en entr\'ee $f_A$ et en sortie le vecteur de d\'eplacements $u$. Des projecteurs $\Pi_A$ et $\Pi_B$ extraient $u_A$ et $u_B$ de $u$.

Q1) Sch\'ema bloc en boucle ouverte :

\begin{tikzpicture}[node distance=2cm, auto]
    \node [input] (input) {};
    \node [block, right of=input] (Fm) {$F_m$};
    \node [block, right of=Fm] (Ft) {$F_t$};
    \node [block, right of=Ft] (F) {$F$};
    \node [output, right of=F, node distance=3cm] (u) {};
    \node [block, below of=F, node distance=2cm] (PiA) {$\Pi_A$};
    \node [block, right of=PiA] (Gt) {$G_t$};
    \node [output, right of=Gt] (theta) {};
    \node [block, above of=F, node distance=2cm] (PiB) {$\Pi_B$};
    \node [output, right of=PiB] (uB) {};

    \path [line] (input) -- node {$i$} (Fm);
    \path [line] (Fm) -- node {$C_m$} (Ft);
    \path [line] (Ft) -- node {$f_A$} (F);
    \path [line] (F) -- node {$u$} (u);
    \path [line] (F) |- (PiA);
    \path [line] (PiA) -- node {$u_A$} (Gt);
    \path [line] (Gt) -- node {$\vartheta$} (theta);
    \path [line] (F) |- (PiB);
    \path [line] (PiB) -- node {$u_B$} (uB);
\end{tikzpicture}

Le moteur est associ\'e \`a un capteur angulaire mesurant $\vartheta = G_t u_A$.

Q2) Relation entre $F_t$ et $G_t$ :

On a : $\theta = G_t u_A \Rightarrow \dot{\theta} = G_t \dot{u_A}$

et $f_A = F_t C_m$

Hypoth\`ese : Conservation de l'\'energie (rendement unitaire)

\[\Rightarrow C_m \dot{\theta} = f_A \dot{u_A}\]

\[\Rightarrow C_m G_t \dot{u_A} = F_t C_m \dot{u_A}\]

\[\Rightarrow G_t = F_t\]

On souhaite $F_t G_t = 1$, c'est-\`a-dire $F_t = 1/G_t$. Cette relation est vraie sous les hypoth\`eses que la transmission est conservative (pas de pertes \'energ\'etiques), sans glissement, et en r\'egime quasi-statique o\`u la puissance est conserv\'ee : $C_m \dot{\vartheta} = f_A \dot{u_A}$.

Q3) Int\'er\^et du jumeau num\'erique pour contr\^oler $u_B$ :

Le syst\`eme industriel ne dispose pas de mesure de $u_B$. Un jumeau num\'erique permet d'estimer $u_B$ \`a partir du mod\`ele et des mesures disponibles (comme $u_A$ via $\vartheta$), afin de le contr\^oler en boucle ferm\'ee sans capteur direct.

Q4) Sch\'ema bloc en boucle ferm\'ee avec mesure de $u_B$ :

On utilise un t\'el\'em\`etre laser avec fonction de transfert $T$ pour mesurer $u_B$.

\begin{tikzpicture}[node distance=2cm, auto]
    \node [input] (input) {};
    \node [sum, right of=input] (sum) {};
    \node [block, right of=sum] (C) {$C$};
    \node [block, right of=C] (Fm) {$F_m$};
    \node [block, right of=Fm] (Ft) {$F_t$};
    \node [block, right of=Ft] (F) {$F$};
    \node [block, right of=F, node distance=3cm] (PiB) {$\Pi_B$};
    \node [output, right of=PiB] (uB) {};
    \node [block, below of=PiB, node distance=3cm] (T) {$T$};

    \path [line] (input) -- node {$e$} (sum);
    \path [line] (sum) -- (C);
    \path [line] (C) -- node {$i$} (Fm);
    \path [line] (Fm) -- node {$C_m$} (Ft);
    \path [line] (Ft) -- node {$f_A$} (F);
    \path [line] (F) -- node {$u$} (PiB);
    \path [line] (PiB) -- node {$u_B$} (uB);
    \path [line] (PiB) |- (T);
    \path [line] (T) -| node [near end] {$-$} (sum);
\end{tikzpicture}

Q5) Sch\'ema bloc du syst\`eme asservi \`a l'aide du jumeau num\'erique (filtre de Kalman) :

On utilise un filtre de Kalman $T_{kalman}$ prenant la commande $C_m$ (apr\`es inversion par $F_0^{-1}$ pour obtenir $f_A$) et la mesure filtr\'ee $\hat{u}_A$ pour estimer $\hat{u}$, puis extraire $\hat{u}_B$ via $\Pi_B$ pour le retour.

\begin{tikzpicture}[node distance=2cm, auto]
    \node [input] (input) {};
    \node [sum, right of=input] (sum) {};
    \node [block, right of=sum] (kctrl) {kctrl};
    \node [block, right of=kctrl] (Fm) {$F_m$};
    \node [block, right of=Fm] (Ft) {$F_t$};
    \node [block, right of=Ft] (F) {$F$};
    \node [block, right of=F, node distance=3cm] (PiB) {$\Pi_B$};
    \node [output, right of=PiB] (uB) {};
    \node [block, below of=F, node distance=2cm] (PiA) {$\Pi_A$};
    \node [output, right of=PiA] (uA) {};
    \node [block, below of=sum, node distance=4cm] (T0) {T0};
    \node [block, right of=T0, node distance=4cm] (Tkalman) {$T_{kalman}$};
    \node [block, right of=Tkalman, node distance=3cm] (Cm_inv) {$C_m^{-1}$};
    \node [output, right of=Cm_inv] (C) {};

    \path [line] (input) -- node {$e$} (sum);
    \path [line] (sum) -- (kctrl);
    \path [line] (kctrl) -- node {$i$} (Fm);
    \path [line] (Fm) -- node {$C_m$} (Ft);
    \path [line] (Ft) -- node {$f_A$} (F);
    \path [line] (F) -- node {$u$} (PiB);
    \path [line] (PiB) -- node {$u_B$} (uB);
    \path [line] (F) -- (PiA);
    \path [line] (PiA) -- node {$u_A$} (uA);
    \path [line] (uA) |- (T0);
    \path [line] (T0) -- node {$\bar{\theta}$} (Tkalman);
    \path [line] (Fm) |- node [near start] {$C_m$} ++(0,-2cm) -| (Tkalman);
    \path [line] (Tkalman) -- (Cm_inv);
    \path [line] (Cm_inv) -- node {$C$} (C);
    \path [line] (Tkalman) -| node [near end] {$\hat{u}_B$} ++(0,4cm) -| node [near end] {$-$} (sum);
\end{tikzpicture}

\subsection{4.1.2. \'Etude du filtre de Kalman}

Q1) Expression de Q :

\[ Q = B q B^T \]

Proc\'edure de construction de Q

Construction de D :

- Construire un second membre $\xi = \begin{bmatrix} 0 \\ \vdots \\ 0 \\ \vdots \\ 0 \end{bmatrix} \leftarrow$ ddt de $u_a$

- utiliser la fonction Newmark1StepMRHS avec les arguments : $(M, C, k, \xi, 0, 0, 0, dt, \beta, \gamma)$ dont les sorties sont les suivantes :

\[ \begin{bmatrix} u, v, a \end{bmatrix} \]

on obtient alors $\beta = \begin{bmatrix} u \\ v \\ a \end{bmatrix}$

Construction de F :

- Construire 2 Matrices :

\[ M_1 = \begin{bmatrix} 1 & 0 \\ 0 & 1 \end{bmatrix} \quad \text{(identit\'e de taille N ddl)} \quad \text{(taille de } u) \]

\[ M_2 = \begin{bmatrix} 0 & 1 \end{bmatrix} \quad \text{(z\'ero de taille N ddl)} \]

Id\'ee: envoyer $\Gamma_3 = \begin{bmatrix} M_1 & M_2 & M_2 \\ M_2 & M_1 & M_2 \\ M_2 & M_2 & M_1 \end{bmatrix}$ dans Newmark,

par morceaux :

\[ U_0 = \begin{bmatrix} M_1 & M_2 & M_2 \end{bmatrix} \]

\[ \alpha_0 = \begin{bmatrix} M_2 & M_1 & M_2 \end{bmatrix} \]

\[ a_0 = \begin{bmatrix} M_2 & M_2 & M_1 \end{bmatrix} \]

on obtient alors :

\[ \begin{bmatrix} u, v, a \end{bmatrix} = \text{Newmark1StepMRHS}(M, C, k, \xi, U_0, \alpha_0, a_0, dt, \beta, \gamma) \]

Construction de $FRF^T$ \`a l'\'etape de pr\'ediction :

une fois $F$ \'evalu\'ee gr\^ace \`a la fonction Newmark \`a l'instant initial, si on a pas d'id\'ee prendre $P$ comme une matrice identit\'e. Elle n'actualisera au cours des pas de temps

Forme de la matrice $H$ :

\[ H = \begin{bmatrix} 1 & 0 \\ 0 & 0 \end{bmatrix} \]

\section{4.2. Travail durant la s\'eance et analyse des r\'esultats}

\subsection{4.2.1. Syst\`eme en boucle ouverte}

Q1) Observation de la position du point B sur la poutre (apr\`es compr\'ehension et lanc\'e du code) :

L'\'evolution de la position du point B est semblable \`a celle du point A (forme sinuso\"idale). Elle poss\`ede un retard d'environ une demi-p\'eriode.

On observe que le mod\`ele reste dans le bon intervalle que la mesure. Mais diff\'erent en amplitude.

Q2) V\'erification de l'invariance de la solution quand on change l'appel de la fonction newmark2 par la boucle en commentaire :

En effet la fonction newmark1stepMRHS effectue un pas de la fonction newmark2 et on obtient la m\^eme solution.

La solution est bien la m\^eme.

Q3) Construction et ajout de la perturbation \`a la force appliqu\'ee en A :

L'ajout de la perturbation en A donne la figure ci-dessous.

On observe que le d\'eplacement au point B est inchang\'e malgr\'e la perturbation en A sur l'effort.

La perturbation a \'et\'e ajout\'ee \`a la force totale au point A. La solution est la m\^eme.

Q4) Cas o\`u la sollicitation est un \'echelon de force :

La solution subit une variation sur une p\'eriode d'environ 7 secondes et se stabilise par la suite. On peut conclure que le sch\'ema num\'erique est fiable dans le temps.

La solution $U_0$ est la m\^eme pour un effort au point A type sinuso\"idale comme type \'echelon. On conclut donc que le mod\`ele num\'erique est stable.

\subsection{4.2.2. Syst\`eme en boucle ferm\'ee avec mesure de $u_B$}

Q1) Calcul analytique de la consigne, de sa d\'eriv\'ee et son int\'egrale :

La consigne est la fonction de Heaviside en t=1s :

\[ H(t - 1) = \begin{cases} 0 & \text{pour } t < 1 \\ 1 & \text{pour } t \ge 1 \end{cases} \]

Sa d\'eriv\'ee est d\'efinie comme suit :

\[ \frac{d}{dt} H(t - 1) = \delta(t - 1) \]

Son int\'egrale est d\'efinie comme suit :

\[ (t)H(t - 1) = \begin{cases} 0 & \text{pour } t < 1 \\ t & \text{pour } t \ge 1 \end{cases} \]

Q2) Impl\'ementation du contr\^oleur PID :

Cas sans bruit de mesure ni perturbation :

On constate que le d\'eplacement est le m\^eme que celui en boucle ouverte.

Cas sans bruit de mesure mais avec perturbation.

Cas avec bruit de mesure et perturbation.

Q3) Le contr\^ole du syst\`eme est-il plus facile ou plus difficile \`a contr\^oler ?

On constate que l'amplitude du d\'eplacement est plus grande pour $nA=nx$ que pour $nA=nx/2$. On peut donc conclure que le syst\`eme est instable. Cette instabilit\'e n'est pas aussi si grande car elle converge vers la consigne apr\`es pratiquement le m\^eme temps de soit 10 it\'erations.

\subsection{4.2.3. Filtre de Kalman}

D\'eplacement et \'ecart types : compatibilit\'e avec ajout de la phase de correction.

\subsection{4.2.4. Syst\`eme en boucle ferm\'ee sans mesure de $u_B$}

Commentaire : le syst\`eme est moins pr\'ecis sur la pr\'ediction de l'\'etat. Il y a plus de bruit. Il est instable.
\newpage
Manipulations Réalisées
1. Construction de la Commande Optimale sans Perturbation :
   - Utilisation de la fonction forthodir pour calculer la commande optimale respectant la consigne de déplacement au bout de la poutre (étape à \(t = 1\) s).
   - Paramètres fixés : \(\mu = 0.1\) (régularisation), \(nopt = 100\), \(\text{stagEps} = 10^{-4}\), intégration Newmark (\(\beta = 0.25\), \(\gamma = 0.5\)).
   - Simulation du système avec la commande optimale pour vérifier le suivi de consigne.
2. Étude de l'Impact de \(\mu\) :
   - Test de \(\mu = 0.01\), \(0.1\), \(1\).
   - Mesure de la vitesse de convergence (itérations, ratio résiduel), qualité de suivi (erreur moyenne/maximale), amplitude de commande.
   - Génération de graphiques comparatifs pour les solutions et résidus.
3. Comparaison avec Cas Perturbé :
   - Calcul de la commande optimale sans perturbation.
   - Simulation avec perturbation ajoutée (bruit sur la force à l'actionneur).
   - Comparaison avec contrôle PID (avec/sans perturbation) via graphiques de déplacement au bout.
Observations
- Sans Perturbation : La commande optimale converge rapidement et assure un suivi quasi-parfait de la consigne, avec une amplitude de commande raisonnable (\(\approx 22\)).
- Impact de \(\mu\) : Peu d'effet observable dans la plage testée ; tous les \(\mu\) convergent au même niveau résiduel, produisant des résultats similaires. \(\mu = 0.1\) est suffisant.
- Avec Perturbation : Le contrôle optimal (ouvert) se dégrade face aux perturbations, tandis que le PID (fermé) maintient une meilleure robustesse. Cela souligne la limitation du contrôle ouvert en environnement incertain.
Interprétations
- Le contrôle optimal par forthodir est efficace pour des scénarios déterministes, offrant une commande anticipative idéale.
- La régularisation \(\mu\) joue un rôle mineur ici, probablement car la convergence est limitée par la tolérance plutôt que par l'ill-conditionnement.
- En présence de perturbations, un contrôle fermé (comme PID) est préférable au contrôle ouvert optimal pour l'adaptation en temps réel.
- Faisabilité du Contrôle Non-Causal (Bonus) : Cette méthode calcule une commande optimale sur \(m > n\) pas futurs, mise à jour tous les \(n\) pas. Elle pourrait améliorer la performance si l'erreur de modèle est non-nulle (e.g., avec Kalman pour estimation d'état), mais le temps de calcul (plusieurs appels à forthodir par fenêtre) la rend peu pratique en temps réel, même avec optimisation. Elle est plus adaptée à des simulations off-line qu'à des applications embarquées.
Ce travail illustre l'intérêt du contrôle optimal pour des trajectoires précises, tout en soulignant l'importance de la robustesse en environnement réel.
\end{document}